% !TeX program = xelatex
% UTF-8. Standalone transcription; compile from this directory or repository root.
% Source: Google Drive / Notebook / ped UTI; filename order matches handwritten pages 1--5.
% pic01.jpg = 20261016_054809709.jpg (1zccuWirA0ggtfaU9Ma5HsGPA3uAuQMwq)
% pic02.jpg = 20261016_054820943.jpg (133xIKTVc8FVfQOXIpAO4TnGYRWOzPBG9)
% pic03.jpg = 20261016_054825482.jpg (1JyDvpEmn7SitGxhlRQOC9xgoTqThlxgz)
% pic04.jpg = 20261016_054831568.jpg (15xhm3sEd4bEROGVVYDSMEz6H8X_ehi9q)
% pic05.jpg = 20261016_054836686.jpg (13UV1On3K9tcPHLzZiiveflQkRYkFdpvl)
% Dual-mode guard: compiles standalone with XeLaTeX, or is \input from the
% XePersian notebook (main.tex) where its local preamble below is skipped.
\newif\ifPedUTIStandalone
\ifdefined\settextfont
  \PedUTIStandalonefalse
\else
  \PedUTIStandalonetrue
\fi
\ifPedUTIStandalone
\documentclass[11pt,a4paper]{article}
\usepackage[margin=20mm]{geometry}
\usepackage{amsmath,amssymb}
\usepackage{graphicx}
\usepackage{array,booktabs}
\usepackage{enumitem}
\usepackage[normalem]{ulem}
\usepackage{xepersian}
\settextfont{Amiri}
\setlatintextfont{DejaVu Serif}
\setlength{\parindent}{0pt}
\setlength{\parskip}{4pt}
\setlength{\emergencystretch}{3em}
\setlist[itemize]{itemsep=3pt,topsep=3pt}
\newcommand{\en}[1]{\lr{#1}}
\newcommand{\yes}{\(\checkmark\)}
\newcommand{\no}{\(\times\)}
\newcommand{\reading}[1]{[خوانش نامطمئن: #1]}
\newcommand{\UTIpagehead}[1]{%
  \section*{تصویر #1 ـ ۵۷۵ ـ \en{UTI}}
  {\small سربرگ چاپی: \en{Subject:}\quad\en{Date: / /} (خالی).
  پابرگ چاپی: \en{YASHA}؛ شمارهٔ دست‌نویس صفحه: #1.}\par}
\begin{document}
\fi

% Helper macros used inside the body must exist in notebook mode too.
\providecommand{\yes}{\(\checkmark\)}
\providecommand{\no}{\(\times\)}
\providecommand{\reading}[1]{[خوانش نامطمئن: #1]}
\providecommand{\UTIpagehead}[1]{%
  \section*{تصویر #1 ـ ۵۷۵ ـ \en{UTI}}
  {\small سربرگ چاپی: \en{Subject:}\quad\en{Date: / /} (خالی).
  پابرگ چاپی: \en{YASHA}؛ شمارهٔ دست‌نویس صفحه: #1.}\par}

\begin{center}
{\Large\bfseries ۵۷۵ ـ \en{UTI}}\par
\end{center}

\UTIpagehead{۱}
\subsection*{درمان \en{UTI} ـ درمان سیستیت حاد}
\begin{itemize}
\item باید سریعاً درمان شود تا از تبدیل آن به پیلونفریت جلوگیری کنیم.
\item اگر علائم شدید و \en{U/A} نشان‌دهندهٔ پیوری باشد، درمان تجربی را شروع می‌کنیم تا نتیجهٔ کشت حاضر شود.
\item اگر علائم خفیف و شک به تشخیص: انتظار برای نتیجهٔ کشت.
\item در \en{Acute Febrile UTI} فرض را بر این می‌گذاریم که پیلونفریت رخ داده است و آن را درمان می‌کنیم.
\item زیر ۱ ماه، شک به \en{UTI} تب‌دار: بستری.
\item بالای ۱ ماه: سرپایی، مگر اندیکاسیون دیگر.
\item انتخاب درمان تجربی بر اساس الگوی مقاومت منطقه‌ای انجام می‌شود.
\item سفالکسین (نسل یک): برای \en{E. coli}.
\item سفیکسیم (نسل سه): در مناطقی که مقاومت \en{E. coli} بالا است.
\item \en{TMP-SMX}: مقاومت به آن در حال بالا رفتن است.
\item نیتروفورانتوئین: صرفاً برای سیستیت؛ به کلیه دوز مناسب نمی‌رسد.
\item در کودک با \en{UTI} تب‌دار نباید استفاده شود.
\item فلوروکینولون خوراکی (سیپروفلوکساسین): سودومونا.
\item آسیب غضروفی به عنوان عارضه.
\item جنتامایسین: حداقل ۱ ماه سن و بستری.
\item نوزادان: آمپی‌سیلین + جنتامایسین یا سفالوسپورین نسل سه.
\item پس از پایان دورهٔ درمان نیاز به \en{U/C} مجدد نیست.
\item \en{U/C} بعد از ۲۴ ساعت از شروع درمان منفی می‌شود.
\item ظرف ۴۸–۷۲ ساعت پس از شروع درمان تب رفع می‌شود.
\end{itemize}

\clearpage
\UTIpagehead{۲}
\begin{itemize}
\item طول درمان: ۳–۵ روز برای سیستیت؛ ۷–۱۰ روز برای پیلونفریت بدون عارضه.
\item اندیکاسیون برای بررسی بیشتر:
\begin{itemize}
\item عدم پاسخ به درمان ظرف ۴۸–۷۲ ساعت.
\item عدم برون‌دهٔ ادراری کافی.
\item افزایش \en{Cr}.
\item لمس توده.
\end{itemize}
\end{itemize}
\subsection*{درمان نفرونیا لوبار حاد، آبسهٔ کلیه، آبسهٔ پری‌نفریک}
\begin{itemize}
\item نفرونیا لوبار حاد با همان \en{AB} پیلونفریت درمان می‌شود، ولی طول دورهٔ درمان ۱۴–۲۱ روز است.
\item آبسه نیاز به درناژ دارد.
\item اول تلاش می‌کنیم از راه پوست؛ اگر نشد، جراحی.
\item آبسهٔ کلیه زیر \en{3 cm} با \en{ABT} به‌تنهایی قابل درمان است.
\item در خصوص آبسه‌های بزرگ‌تر، اگر بیمار علامت وخیمی نداشته باشد، می‌توان ۴۸ ساعت \en{ABT} را امتحان کرد؛ اگر جواب نداد، درناژ.
\item آبسه نیازمند ۱۰–۱۴ روز درمان \en{AB} وریدی است و به دنبال آن ۲–۴ هفته خوراکی (اگر کشت نداشتیم: \en{E. coli} و استاف اورئوس).
\end{itemize}

\clearpage
\UTIpagehead{۳}
\subsection*{اندیکاسیون‌های تصویربرداری برای شیرخوار با \en{UTI} تب‌دار}
\begin{itemize}
\item همه سونوگرافی انجام بدهند.
\item \en{VCUG} اگر:
\begin{itemize}
\item سونوگرافی یافتهٔ غیرطبیعی داشت.
\item تکرار اپیزود \en{UTI} تب‌دار.
\end{itemize}
\end{itemize}
\subsection*{اندیکاسیون‌های تصویربرداری برای کودکان با \en{UTI}}
{\small سه ستون علامت در دست‌نویس عنوان ندارند. ترتیب آن‌ها در جدول‌های زیر از راست به چپ، مطابق تصویر حفظ شده است؛ \yes{} و \no{} عین علامت‌های تصویرند.}

\textbf{زیر ۶ ماه}
\begin{itemize}
\item سونوگرافی در فاز حاد:
\begin{itemize}
\item پاسخ ظرف ۴۸ ساعت: \no.
\item عفونت آتیپیک: \yes{} (شدیداً \en{ill}، کاهش ادرار، سپسیس، \en{\(Cr\uparrow\)}).
\item عفونت راجعه: \yes.
\end{itemize}
\end{itemize}
\begin{center}
\begin{tabular}{p{8cm}ccc}
\toprule
بررسی & ۱ & ۲ & ۳ \\
\midrule
سونوگرافی ظرف ۶ هفتهٔ بعد & \yes & \no & \no \\
\en{VCUG} (اگر سونو نرمال) & \yes & \yes & \yes \\
\bottomrule
\end{tabular}
\end{center}
{\small فلش به علامت اولِ سونوگرافی ظرف ۶ هفتهٔ بعد در تصویر رسم شده است.}

\textbf{۶ ماه تا ۳ سال}
\begin{center}
\begin{tabular}{p{8cm}ccc}
\toprule
بررسی & ۱ & ۲ & ۳ \\
\midrule
سونو در فاز حاد & \no & \yes & \no \\
سونو ظرف ۶ هفته & \no & \no & \yes \\
\bottomrule
\end{tabular}
\end{center}
\en{VCUG}: \no{}؛ \en{not routine}.
{\small فلش به علامت اولِ سونو در فاز حاد در تصویر رسم شده است.}

\textbf{بالاتر از ۳ سال}
\begin{center}
\begin{tabular}{p{8cm}ccc}
\toprule
بررسی & ۱ & ۲ & ۳ \\
\midrule
سونو در فاز حاد & \no & \yes & \no \\
سونو ظرف ۶ هفتهٔ بعد & \no & \no & \yes \\
\en{VCUG} & \no & \no & \no \\
\bottomrule
\end{tabular}
\end{center}

\clearpage
\UTIpagehead{۴}
\subsection*{پروفیلاکسی}
کدام بیماری‌ها پس از اپیزود اول \en{UTI} تب‌دار از \en{AB} پروفیلاکسی سود می‌برند؟ البته در برخی اندیکاسیون‌ها:
\begin{itemize}
\item مثانهٔ نوروژنیک.
\item انسداد و اتساع مجاری ادراری.
\item \en{Severe VUR}.
\item سنگ‌های ادراری.
\end{itemize}
پروفیلاکسی روتین برای \en{VUR} توصیه نمی‌شود، خصوصاً گریدهای پایین.

\subsection*{تصویربرداری در \en{UTI}}
{\small واژهٔ «تشخیص» در عنوان دست‌نویس خط خورده و «تصویربرداری در» بالای آن نوشته شده است.}
\begin{itemize}
\item نیازمند تصویربرداری نیست.
\item اگر به درمان جواب ندهد: تصویربرداری.
\item سونوگرافی خط اول تصویربرداری‌ها است.
\item \en{CT} برای تشخیص لوبار نفرونیا حساس‌تر است.
\item پس از کنتراست \en{low density} می‌شود.
\end{itemize}

\subsection*{تشخیص \en{UTI}}
\begin{itemize}
\item ریسک فاکتور \en{UTI} در شیرخواران زیر ۲ سال:
\begin{itemize}
\item سن زیر ۱ سال، دختر بودن، پسر ختنه‌نشده، تب \en{\(\geq39\)}، تب بیش از ۲ روز، نبود سورس دیگر عفونت.
\end{itemize}
\item اگر ریسک \en{UTI} \en{\(\geq2\%\)}: \en{U/A} و \en{U/C}.
\item \en{U/A}: پیوری + \en{\(\geq50{,}000\;CFU\)}؛ یا \en{\(\geq10{,}000\;CFU\)} در زمینهٔ بالینی مناسب.
\end{itemize}

\clearpage
\UTIpagehead{۵}
\begin{itemize}
\item زیر ۲ سال باید با سونداژ نمونه گرفته شود.
\item \en{bag} ارزش ندارد، مگر منفی باشد.
\item سوپراپوبیک جایگاهی ندارد.
\item پیوری یعنی \en{\(\geq5\;WBC\)} در نمونهٔ سانتریفیوژشده؛ \en{\(\geq10\;WBC\)} در نمونهٔ \en{U/A}؛ یا لکوسیت استراز + در دیپ‌استیک.
\item نیتریت در \en{U/A} از تشخیص \en{UTI} حمایت می‌کند.
\item در انتروکوک منفی است.
\item هماچوری میکروسکوپیک در سیستیت حاد شایع است.
\item هماچوری تنها به معنای \en{UTI} نیست.
\item \en{WBC Cast} در سدی‌مان ادراری مطرح‌کنندهٔ درگیری کلیه است.
\end{itemize}

\subsection*{پیوری استریل}
تعریف: کشت منفی + لکوسیت.

\subsection*{در درگیری حاد کلیه}
لکوسیتوز، نوتروفیلی، \en{\(ESR\uparrow\)}، \en{\(\uparrow\)} پروکلسی‌تونین، \en{\(CRP\uparrow\)}.

\subsection*{باکتریمی}
در ۲–۳٪ موارد پیلونفریت رخ می‌دهد، خصوصاً در زیر ۶۰ روز.
\begin{itemize}
\item کودک \en{ill} یا زیر ۶۰ روز: \en{B/C}.
\end{itemize}

\clearpage
\section*{تصاویر اصلی}
\IfFileExists{images/pic01.jpg}{\def\imagebase{images/}}{%
\def\imagebase{pediatrics/ped-infectious diseases/UTI/images/}}
\newcommand{\sourceimage}[2]{%
\begin{center}
\textbf{تصویر #1}\par
\includegraphics[width=\linewidth,height=0.88\textheight,keepaspectratio]{\imagebase#2}
\end{center}}
\sourceimage{۱}{pic01.jpg}
\clearpage
\sourceimage{۲}{pic02.jpg}
\clearpage
\sourceimage{۳}{pic03.jpg}
\clearpage
\sourceimage{۴}{pic04.jpg}
\clearpage
\sourceimage{۵}{pic05.jpg}
\ifPedUTIStandalone
\end{document}
\fi
